\section{Successive Schmidt trimming}
\label{sec:lvv_successive_trimming}

The bond-trimming step of Landau, Vazirani and
Vidick~\cite[Claim~2.10, source lines 549--559]{Landau2015Polynomial} trims every
member of a viable set successively across all internal cuts. This section proves
Theorem~\ref{thm:lvv_size_bond_trimmed_family}: after size trimming and bond
trimming, a family chosen before the ground state is $1/2$-viable, with polynomial
cardinality and polynomial bond dimension. The retained dimension is quadratic in the
chain length, in place of the printed linear one~\cite{gap:lvv13_trimming_radius}.

\begin{lemma}[Numbering the sites of a left block]
    \label{lem:lvv_left_coordinates}
    \lean{TNLean.MPS.AreaLaw.leftConfigOfFinEquiv,
        TNLean.MPS.AreaLaw.rank_cutMatrix_leftConfigOfFin}
    \leanok
    \uses{def:lvv_viable_mps}
    For $i\le L$, configurations of the first $i$ sites of an $L$-site chain
    are canonically identified with configurations of an $i$-site chain.
    This identification preserves the rank of the coefficient matrix at
    every cut after $j$ sites, including the endpoint cuts.
\end{lemma}
\begin{proof}
    \leanok
    The identification sends each site to the site with the same index.
    At a cut, it permutes the row and column configurations independently;
    such permutations preserve matrix rank.
\end{proof}

\begin{definition}[Successive Schmidt trimming]
    \label{def:lvv_successive_trimming}
    \lean{TNLean.MPS.AreaLaw.chainTrimmingProjection,
        TNLean.MPS.AreaLaw.isTrimmingProjection_chainTrimmingProjection,
        TNLean.MPS.AreaLaw.successiveSchmidtTrim,
        TNLean.MPS.AreaLaw.columnwiseSchmidtTrim}
    \leanok
    \uses{def:chain_model}
    Fix a retained dimension $m\in\N$ and a chain vector $\psi$.
    Set $\psi_0=\psi$. At step $j+1$, let $P_{j+1}$ be an orthogonal
    projection onto at most $m$ leading left Schmidt vectors of $\psi_j$
    across the cut after $j+1$ sites, and set
    $\psi_{j+1}=(P_{j+1}\otimes I)\psi_j$.
    Choices within a repeated Schmidt value are allowed. For a vector
    $u=\sum_a u_a\otimes|a\rangle$ with a finite environment, independent
    column trimming replaces each $u_a$ by its successive trimming through
    all internal bonds. This is the individual-vector operation in the
    bond-trimming step of~\cite[source lines 549--559]{Landau2015Polynomial}.
\end{definition}

\begin{lemma}[Norm and bond dimensions after successive trimming]
    \label{lem:lvv_successive_trimming_rank}
    \lean{TNLean.MPS.AreaLaw.norm_successiveSchmidtTrim_le,
        TNLean.MPS.AreaLaw.rank_cutMatrix_successiveSchmidtTrim_le,
        TNLean.MPS.AreaLaw.rank_cutMatrix_successiveSchmidtTrim_le_of_le,
        TNLean.MPS.AreaLaw.isOpenMPS_successiveSchmidtTrim,
        TNLean.MPS.AreaLaw.norm_columnwiseSchmidtTrim_le}
    \leanok
    \uses{def:lvv_successive_trimming,def:chain_open_mps}
    Every successive trimming satisfies $\|\psi_j\|\le\|\psi\|$ and
    $\operatorname{rank}M_k(\psi_j)\le\operatorname{rank}M_k(\psi)$ at
    every cut $k$. If $0<k\le j$, then
    $\operatorname{rank}M_k(\psi_j)\le m$. Consequently, trimming the
    $L-1$ internal bonds gives an open matrix product state of bond
    dimension at most $m$. Independent column trimming does not increase
    the joint norm $\|u\|$.
\end{lemma}
\begin{proof}
    \leanok
    \uses{def:chain_trim,thm:chain_trim_cut_ranks}
    Each trimming is an orthogonal projection and therefore contracts
    the norm. A Schmidt trimming does not increase any cut rank. At the
    processed cut its rank is at most $m$, and the later trimmings preserve
    this bound. The characterization of open matrix product states by
    internal cut ranks gives the representation. For column trimming,
    sum $\|\widetilde u_a\|^2\le\|u_a\|^2$ over the environment index.
\end{proof}

\begin{lemma}[Accumulated overlap loss]
    \label{lem:lvv_successive_trimming_overlap}
    \lean{TNLean.MPS.AreaLaw.norm_inner_sub_successiveSchmidtTrim_le}
    \leanok
    \uses{def:lvv_successive_trimming}
    Let $v$ be a chain vector whose coefficient matrix has rank at most
    $D$ at each of the first $j$ cuts. Then
    \begin{align}
        |\langle\psi-\psi_j,v\rangle|
         & \le j\|\psi\|\|v\|\sqrt{\frac{D}{m+1}}.
    \end{align}
    Thus the overlap loss is at most $\varepsilon$ for unit vectors when
    $m+1\ge j^2D/\varepsilon^2$. This uses the corrected single-cut
    radius of~\cite{gap:lvv13_trimming_radius}.
\end{lemma}
\begin{proof}
    \leanok
    \uses{thm:chain_trim_error,lem:lvv_successive_trimming_rank}
    Write $\psi-\psi_j=\sum_{k=0}^{j-1}(\psi_k-\psi_{k+1})$.
    The corrected single-cut estimate bounds the pairing of the $k$th
    summand with $v$ by $\|\psi_k\|\|v\|\sqrt{D/(m+1)}$.
    Use $\|\psi_k\|\le\|\psi\|$ and the triangle inequality.
\end{proof}

\begin{lemma}[Overlap loss for independently trimmed columns]
    \label{lem:lvv_columnwise_trimming_overlap}
    \lean{TNLean.MPS.AreaLaw.norm_inner_sub_columnwiseSchmidtTrim_le,
        TNLean.MPS.AreaLaw.norm_inner_sub_leftColumnwiseSchmidtTrim_le}
    \leanok
    \uses{def:lvv_successive_trimming}
    Let $u=\sum_a u_a\otimes|a\rangle$ and
    $v=\sum_a v_a\otimes|a\rangle$ be vectors on an $L$-site chain and a
    finite environment with its orthonormal coordinate basis. Suppose that
    every $v_a$ has coefficient-matrix rank at most $D$ at every internal
    cut. If $\widetilde u$ is obtained by independently trimming each
    $u_a$ to bond dimension $m$, then
    \begin{align}
        |\langle u-\widetilde u,v\rangle|
         & \le (L-1)\|u\|\|v\|\sqrt{\frac{D}{m+1}}.
    \end{align}
\end{lemma}
\begin{proof}
    \leanok
    \uses{lem:lvv_successive_trimming_overlap}
    The inner product is the sum of the inner products of the columns.
    Lemma~\ref{lem:lvv_successive_trimming_overlap} bounds each summand by
    $(L-1)\|u_a\|\|v_a\|\sqrt{D/(m+1)}$. Cauchy--Schwarz gives
    $\sum_a\|u_a\|\|v_a\|\le\|u\|\|v\|$, since the environment coordinates
    are orthonormal.
\end{proof}

\begin{definition}[Bond trimming of a finite left family]
    \label{def:lvv_bond_trimmed_family}
    \lean{TNLean.MPS.AreaLaw.leftSchmidtTrim,
        TNLean.MPS.AreaLaw.leftSchmidtTrim_leftConfigOfFin,
        TNLean.MPS.AreaLaw.bondTrimSet,
        TNLean.MPS.AreaLaw.leftColumnwiseSchmidtTrim}
    \leanok
    \uses{def:lvv_successive_trimming,def:lvv_viable_mps}
    Fix $i\le L$ and a finite family $S$ of vectors on the first $i$ sites.
    Write $\mathcal T_m u$ for successive Schmidt trimming of $u$ through
    its $i-1$ internal bonds, retaining dimension at most $m$ at each step.
    The trimmed family is $\mathcal T_m(S)=\{\mathcal T_m u:u\in S\}$.
    For $u=\sum_a u_a\otimes|a\rangle$ on the left block and an environment,
    independent column trimming is
    $\widetilde u=\sum_a\mathcal T_m u_a\otimes|a\rangle$.
    These operations depend only on the input vectors and $m$.
\end{definition}

\begin{lemma}[Cardinality and bond dimensions of the trimmed family]
    \label{lem:lvv_bond_trimmed_family_bounds}
    \lean{TNLean.MPS.AreaLaw.isLeftOpenMPS_leftSchmidtTrim,
        TNLean.MPS.AreaLaw.norm_leftSchmidtTrim_le,
        TNLean.MPS.AreaLaw.card_bondTrimSet_le,
        TNLean.MPS.AreaLaw.isLeftOpenMPS_of_mem_bondTrimSet,
        TNLean.MPS.AreaLaw.norm_leftColumnwiseSchmidtTrim_le,
        TNLean.MPS.AreaLaw.leftColumnwiseSchmidtTrim_column_mem}
    \leanok
    \uses{def:lvv_bond_trimmed_family}
    The trimmed family has cardinality at most $|S|$, and each of its
    vectors is an open matrix product state with bond dimension at most
    $\max\{1,m\}$. One has $\|\mathcal T_m u\|\le\|u\|$ and
    $\|\widetilde u\|\le\|u\|$. If every $u_a$ belongs to $S$, every column
    of $\widetilde u$ belongs to $\mathcal T_m(S)$.
\end{lemma}
\begin{proof}
    \leanok
    \uses{lem:lvv_left_coordinates,lem:lvv_successive_trimming_rank}
    Number the first $i$ sites in order and apply successive trimming.
    Lemma~\ref{lem:lvv_successive_trimming_rank} gives the norm and rank
    bounds. An image of a finite set has at most the original cardinality.
    The column statement is the defining image membership.
\end{proof}

\begin{lemma}[Physical error of independent column trimming]
    \label{lem:lvv_physical_column_trimming}
    \lean{TNLean.MPS.AreaLaw.rank_cutMatrix_rightApply_blockSplit_column_le,
        TNLean.MPS.AreaLaw.leftColumnwiseSchmidtTrim_mem_leftSupportSubmodule,
        TNLean.MPS.AreaLaw.norm_inner_sub_leftColumnwiseSchmidtTrim_rightApply_le}
    \leanok
    \uses{def:lvv_bond_trimmed_family}
    Let $i\le L$, let $W:\mathcal E\to\mathcal H_R$ be an isometry from a
    finite environment, and let $u\in\mathcal H_A\otimes\mathcal E$, where
    $\mathcal H_A$ is the first $i$ sites. Let $\widetilde u$ be its
    independently trimmed columns, and write
    $\phi=(I\otimes W)u$, $w=(I\otimes W)\widetilde u$ as chain vectors.
    If $v$ has cut rank at most $D$ at every cut strictly inside
    $\mathcal H_A$, then
    \begin{align}
        |\langle\phi-w,v\rangle|
         & \le(i-1)\|u\|\|v\|\sqrt{\frac{D}{m+1}}.
    \end{align}
    If every column of $u$ belongs to $S$, then $w$ belongs to
    $\operatorname{Span}(\mathcal T_m(S))\otimes\mathcal H_R$.
\end{lemma}
\begin{proof}
    \leanok
    \uses{lem:lvv_columnwise_trimming_overlap,
        lem:lvv_bond_trimmed_family_bounds,thm:lvv_left_columns_right_extension,
        thm:alvv_collective_trim_coordinates}
    Put $v'=(I\otimes W^*)v$. Contracting the environment and then taking a
    column cannot increase an internal cut rank, so every column of $v'$
    has cut rank at most $D$. Since $W$ is an isometry, $\|v'\|\le\|v\|$.
    Apply Lemma~\ref{lem:lvv_columnwise_trimming_overlap} to
    $\langle\phi-w,v\rangle=\langle u-\widetilde u,v'\rangle$.
    The support assertion follows because every column of $\widetilde u$
    belongs to $\mathcal T_m(S)$, and right extension preserves its span.
\end{proof}

\begin{lemma}[Viability after independent bond trimming]
    \label{lem:lvv_bond_trimmed_family_viable}
    \lean{TNLean.MPS.AreaLaw.bondTrimmingRadius,
        TNLean.MPS.AreaLaw.mul_sqrt_bondTrimmingRadius_le,
        TNLean.MPS.AreaLaw.isViableMPS_bondTrimSet_of_columns}
    \leanok
    \uses{def:lvv_bond_trimmed_family,def:lvv_viable_mps}
    Let $i\le L$, let $\Gamma,v$ be unit chain vectors, and suppose that
    $|\langle\Gamma,v\rangle|\ge47/48$. Suppose that $v$ has cut rank at
    most $D$ at every cut strictly inside the first $i$ sites. Let $u$ be
    a unit vector on those sites and a finite environment, let $W$ be an
    isometry into the remaining sites, and suppose that every left column
    of $u$ belongs to the finite set $S$ and
    $|\langle\Gamma,(I\otimes W)u\rangle|\ge11/12$.
    With $m=(48L)^2D$, the independently trimmed family $\mathcal T_m(S)$
    is $(i,|S|,m,1/2)$-viable for $\Gamma$.
    The quadratic radius replaces the printed linear radius
    of~\cite[Claim~2.10, source lines 549--559]{Landau2015Polynomial}, by
    the corrected single-cut estimate~\cite{gap:lvv13_trimming_radius}.
\end{lemma}
\begin{proof}
    \leanok
    \uses{lem:lvv_physical_column_trimming,lem:lvv_bond_trimmed_family_bounds}
    For $\phi=(I\otimes W)u$, the overlap inequality through the common
    unit vector $\Gamma$ gives $|\langle\phi,v\rangle|\ge19/24$.
    The radius satisfies
    $L\sqrt{D/(m+1)}\le1/48$. Thus the independently trimmed extension
    $w$ has $|\langle w,v\rangle|\ge37/48$ and $0<\|w\|\le1$.
    Normalize $z=w/\|w\|$. Then $|\langle z,v\rangle|\ge37/48$, and
    $|\langle\Gamma,z\rangle|\ge1-2(1/48+11/48)=1/2$.
    The vector $z$ is supported on the trimmed family; its cardinality
    and bond bounds follow from
    Lemma~\ref{lem:lvv_bond_trimmed_family_bounds}.
\end{proof}

\begin{lemma}[A polynomial-rank reference for trimming]
    \label{lem:lvv_bond_trimming_reference}
    \lean{TNLean.MPS.AreaLaw.lvvReferenceBond,
        TNLean.MPS.AreaLaw.exists_bond_trimming_reference}
    \leanok
    \uses{def:chain_model,def:chain_open_mps}
    Let $L\ge1$, $q\ge1$ and $\Delta>0$. For every normalized gapped ground
    vector $\Gamma$ of a range-one chain Hamiltonian with term norms at
    most one, there is a unit open matrix product state $v$ such that
    $|\langle\Gamma,v\rangle|\ge47/48$, with bond dimension at most
    $r(L)=D(48L)$, where $D$ is the bound of
    Theorem~\ref{thm:chain_mps_approximation} at $R=J=1$.
    The function $r(L)$ is fixed independently of $\Gamma$ and has a
    polynomial bound in $L$.
\end{lemma}
\begin{proof}
    \leanok
    \uses{thm:chain_mps_approximation_error}
    Apply the arbitrary-accuracy chain approximation with
    $\varepsilon=1/(2304L)$. Its bond bound is $D(48L)$, because
    $2Lc_0/\varepsilon=2(48L)^2c_0$. Cauchy--Schwarz gives
    $|\langle\Gamma,v\rangle|\ge1-\|v-\Gamma\|\ge47/48$.
    The polynomial estimate for $D$ remains polynomial after replacing
    $L$ by $48L$.
\end{proof}

\begin{lemma}[Polynomial growth of the trimming radius]
    \label{lem:lvv_bond_trimming_radius_polynomial}
    \lean{TNLean.MPS.AreaLaw.lvvReferenceBond_le_mul_rpow,
        TNLean.MPS.AreaLaw.bondTrimmingRadius_lvvReferenceBond_le_mul_rpow,
        TNLean.MPS.AreaLaw.exists_bondTrimmingRadius_lvvReferenceBond_polynomial}
    \leanok
    \uses{lem:lvv_bond_trimming_reference}
    Fix $q$ and $\Delta$, and put
    \[
        s=\sqrt{(1+\Delta^{-1})B_1(q)},\qquad
        K=\exp(C_{\mathrm{int}}(q,\Delta))(2c_0)^{64s},
    \]
    where $C_{\mathrm{int}}$, $B_1$ and $c_0=\sqrt{2e^{e/2}}$ are the
    entropy bound, budget constant and tail prefactor from the chain
    approximation theorem. For $L\ge1$,
    \[
        r(L)\le K(48L)^{128s},\qquad
        (48L)^2r(L)\le K(48L)^{128s+2}.
    \]
    In particular, there are $C_1\ge0$ and $p_1\in\mathbb N$, depending
    only on $q,\Delta$, such that $(48L)^2r(L)\le C_1L^{p_1}$ for
    every $L\ge1$.
\end{lemma}
\begin{proof}
    \leanok
    \uses{thm:chain_mps_approximation}
    Apply the polynomial bond bound at length $48L$, and multiply by
    $(48L)^2$. Take $p_1=\lceil128s+2\rceil$ and $C_1=K48^{p_1}$.
    Since $48L\ge1$, increasing the exponent preserves the inequality.
\end{proof}

\begin{theorem}[Size and bond trimming before the ground vector is supplied]
    \label{thm:lvv_size_bond_trimmed_family}
    \lean{TNLean.MPS.AreaLaw.ChainHamiltonian.exists_size_bond_trimmed_family_polynomial}
    \leanok
    \uses{def:chain_model,def:lvv_viable_mps}
    Let $L\ge1$, $q\ge2$, $i\le L$, and $\Delta>0$. Fix a range-one chain
    Hamiltonian with positive semidefinite terms of norm at most one,
    and a finite left family $S$ whose members have bond dimension at
    most $\max\{1,b\}$. Use $c,D,p,A,C$ from
    Theorem~\ref{thm:lvv_size_trimmed_family_polynomial}, and put
    $m=(48L)^2r(L)$, with $r(L)$ from
    Lemma~\ref{lem:lvv_bond_trimming_reference}. This bond radius is
    polynomial in $L$ by Lemma~\ref{lem:lvv_bond_trimming_radius_polynomial}.
    There is a finite left family $T$, chosen before every ground-state
    energy $E_0$ and ground vector $\Gamma$, with $|T|\le CL^p$ and every
    member of bond dimension at most $\max\{1,m\}$, such that
    $(i,c/L)$-viability of $S$ for a normalized ground vector of energy
    $E_0$ and gap at least $\Delta$ implies $(i,1/2)$-viability of $T$.
    This is~\cite[Claim~2.10, source lines 549--559]{Landau2015Polynomial}
    for the size-trimming family of
    Theorem~\ref{thm:lvv_size_trimmed_family_polynomial}, with the bond
    $m=(48L)^2r(L)$ in place of the printed $48L\,r(L)$.
\end{theorem}
\begin{proof}
    \leanok
    \uses{thm:lvv_size_trimmed_family_polynomial,
        lem:lvv_bond_trimming_reference,lem:lvv_bond_trimmed_family_viable}
    Choose the size-trimming family $T_0$ before $E_0,\Gamma$, retaining its
    orthonormal right-column witness. Define $T=\mathcal T_m(T_0)$.
    Independent trimming preserves the cardinality bound and gives the
    stated bond bound. For each admissible ground pair, the size step
    supplies a unit $u$, an isometry $W$, all columns of $u$ in $T_0$, and
    overlap at least $11/12$ for $(I\otimes W)u$.
    The reference theorem supplies $v$ of cut rank at most $r(L)$ and
    overlap at least $47/48$. Apply
    Lemma~\ref{lem:lvv_bond_trimmed_family_viable}.
\end{proof}
